\section{Methods}

To derive our dataset, we largely follow the methods described in \citet{meinshausen_historical_2017}.
However, we have made changes in a number of places
so we describe the method in full,
while also highlighting places of clear divergence.

\subsection{General approach}

In line with \citet{meinshausen_historical_2017},
the general approach is to first create a dataset with monthly temporal resolution
and 15\textdegree latitudinal spatial resolution
(but no longitudinal resolution).
We refer to this as our `native resolution'.
From our native resolution, we can then create higher and lower resolution data products (see [TODO: forward refernence]).

To create our native resolution dataset, we sum three components:

\begin{enumerate}
  \item Global-, annual-mean concentrations
  \item Seasonal variations
  \item Latitudinal gradient
\end{enumerate}

In order to take this sum,
all three components must already have the same time dimension
(monthly timesteps in our case).
This introduces the first subtlety: while our approach is based on a timeseries of global-, annual-mean concentrations,
this timeseries has to be interpolated to monthly timesteps
in order to avoid the change from year to year resulting in step changes between December of one year and January of the next.
The same interpolation requirement applies to the latitudinal gradient
(which is initially calculated on a yearly timestep, see [TODO forward ref]).

In addition to the time dimension consistency requirement, seasonality can vary by latitude.
As a result, the latitudinal variation of seasonality must be calculated
(i.e. seasonality must have a "latitude" dimension)
before taking the sum.

Thus, in practice, we calculate each dataset according to equation \ref{eq:native-resolution-sum}.

\begin{equation}
    \label{eq:native-resolution-sum}
    gn(y, m, l) = gm(y, m) + lg(y, m, l) + s(y, m, l)
\end{equation}
where $gn$ is the native resolution dataset value
(`gn' is the `grid label' used for such output on ESGF, so we re-use it here for convenience and consistency with the published output, see [TODO forward ref]),
$y$ is the year in which the value applies,
$m$ is the month in which the value applies,
$l$ is the latitude at which the value applies,
$gm$ is the global-, annual-mean derived value (and is broadcast across all latitudes),
$lg$ is the latitudinal gradient contribution
and $s$ is the seasonality contribution.

The key choices are about how we arrive at $gm$, $lg$ and $s$ for each greenhouse gas species in our dataset.
We describe the different methods in the following sections.

\subsection{Calculating the general components}

\subsubsection{CO_2}

CO_2 has the longest observational record of all the greenhouse gases included in our dataset.
It is also the largest driver of climate change to date [TODO citation], hence has the most available observations.
Finally, atmospheric CO_2 concentrations are driven by the carbon cycle,
which has led to the introduction of extra complexity in the calculation of the input components,
with the aim of better representing the variations seen in the observational record.

We begin by collecting observations of surface CO_2 concentrations from the NOAA observational networks.
We bin these observations into 15\textdegree latitudinal, 60\textdegree longitudinal and monthly temporal bins.
Within each bin, we firstly take the mean of all measurements that come from the same network
(where, for NOAA, we treat the in-situ and surface flask measurements as separate networks),
station and measurement method.
[TODO figure out what different measurement methods are]
As a result, each network, station and measurement method gets an equal weight in each bin,
irrespective of how many measurements that network-station-measurement made in the bin.
Then, we take the average over network, station and measurement method
to get a single value in each bin.

The surface observational network does not cover the entire globe [TODO: figure],
hence the next step is to interpolate the available measurements to get complete spatial coverage.

- decision about which timesteps we can interpolate for
- use max. number of months for which we can create continuous timeseries
- do stupid linear interpolation [TODO check what scipy griddata does exactly, note that we use linear interpolation and 'wrap around' to ensure we don't get unwanted edge effects]


\begin{enumerate}
  \item NOAA
  \item Seasonality (year, month, latitude)
  \item Latitudinal gradient (year, month, latitude)
\end{enumerate}

For all GHGs, we need three things:

\begin{enumerate}
  \item Global-mean concentration with monthly temporal resolution (explain somewhere why this doesn't work with annual-mean resolution because of the jumps you get from December to January)
  \item Seasonality (year, month, latitude)
  \item Latitudinal gradient (year, month, latitude)
\end{enumerate}

The trick is how we get these for different gases.
Breaks down into six cases:

\begin{enumerate}
    \item CO_2
    \item CH_4
    \item N_2O
    \item SF_6-like
    \item C_4F_10-like
    \item C_8F_18-like
\end{enumerate}

CO_2:

\begin{enumerate}
    \item Bin obs
    \item Interpolate/extrapolate to complete spatial coverage
    \item Do PCA to get global-, annual-mean, seasonality and latitudinal gradient over the observational era with sufficient data
    \item Extend latitudinal gradient PCs back in time using fossil emissions
    \item Extend global-mean back in time (depends on the previous step)
    \item Extend seasonality back in time using our funny temperature-concentration composite
    \item Put it all back together to get the pieces for gridding
    \begin{itemize}
        \item Get annual-mean monthly
        \item Calculate seasonality over time (i.e. base seasonality plus seasonality change)
        \item Interpolate latitudinal gradient PCs to monthly
    \end{itemize}
\end{enumerate}

CH_4:

\begin{enumerate}
    \item Bin obs
    \item Interpolate/extrapolate to complete spatial coverage
    \item Do PCA to get global-, annual-mean, seasonality and latitudinal gradient over the observational era with sufficient data
    \item Extend latitudinal gradient PCs back in time using fossil emissions
    \item Extend global-mean back in time (depends on the previous step)
    \item Put it all back together to get the pieces for gridding
    \begin{itemize}
        \item Get annual-mean monthly
        \item Calculate seasonality over time (i.e. multiply base seasonality by global-mean values)
        \item Interpolate latitudinal gradient PCs to monthly
    \end{itemize}
\end{enumerate}

N_2O:

\begin{enumerate}
    \item Bin obs
    \item Interpolate/extrapolate to complete spatial coverage
    \item Do PCA to get global-, annual-mean, seasonality and latitudinal gradient over the observational era with sufficient data
    \item Extend latitudinal gradient PCs back in time using total emissions
    \item Extend global-mean back in time (depends on the previous step)
    \item Put it all back together to get the pieces for gridding
    \begin{itemize}
        \item Get annual-mean monthly
        \item Calculate seasonality over time (i.e. multiply base seasonality by global-mean values)
        \item Interpolate latitudinal gradient PCs to monthly
    \end{itemize}
\end{enumerate}

SF_6:

\begin{enumerate}
    \item Bin obs
    \item Interpolate/extrapolate to complete spatial coverage
    \item Do PCA to get global-, annual-mean, seasonality and latitudinal gradient over the observational era with sufficient data
    \item Extend latitudinal gradient PCs back in time using total emissions
    \item Extend global-mean back in time using pre-industrial value plus interpolation
    \item Put it all back together to get the pieces for gridding
    \begin{itemize}
        \item Get annual-mean monthly
        \item Calculate seasonality over time (i.e. multiply base seasonality by global-mean values)
        \item Interpolate latitudinal gradient PCs to monthly
    \end{itemize}
\end{enumerate}

C_4F_10:

\begin{enumerate}
    \item Get latitudinal gradient
    \item Put it all back together to get the pieces for gridding
    \begin{itemize}
        \item Get annual-mean monthly
        \item Assume zero seasonality
        \item Interpolate latitudinal gradient PCs to monthly
    \end{itemize}
\end{enumerate}

C_8F_18:

\begin{enumerate}
    \item No new data, so just use CMIP6 pieces (historical plus ssp245, minor gas so different choices make negligible difference [TODO quantify])
    \item Put it all back together to get the pieces for gridding
    \begin{itemize}
        \item Get annual-mean monthly
        \item Assume zero seasonality
        \item Interpolate latitudinal gradient PCs to monthly
    \end{itemize}
\end{enumerate}

\subsection{Details}

Some methods details here CO_2.

